\subsection{Thiết kế API cho phân hệ RAG}

API RAG được sử dụng khi câu hỏi của người dùng không thể trả lời tốt bằng rule, response tĩnh hoặc dữ liệu có cấu trúc. Đây thường là các câu hỏi dài, câu hỏi mở hoặc câu hỏi cần dựa vào tài liệu tuyển sinh.

Nhóm API RAG đề xuất gồm:

\begin{longtable}{|p{3cm}|p{5cm}|p{7cm}|}
\hline
\textbf{Phương thức} & \textbf{Endpoint đề xuất} & \textbf{Mô tả} \\
\hline
\texttt{POST} & \texttt{/v1/rag/query} & Gửi câu hỏi đến phân hệ RAG và nhận câu trả lời kèm nguồn. \\
\hline
\texttt{POST} & \texttt{/v1/rag/documents} & Thêm tài liệu mới vào kho tri thức. \\
\hline
\texttt{GET} & \texttt{/v1/rag/documents} & Lấy danh sách tài liệu đã ingest. \\
\hline
\texttt{GET} & \texttt{/v1/rag/documents/\{id\}} & Xem chi tiết tài liệu. \\
\hline
\texttt{POST} & \texttt{/v1/rag/documents/\{id\}/ingest} & Thực hiện ingest tài liệu vào kho tri thức. \\
\hline
\texttt{POST} & \texttt{/v1/rag/documents/\{id\}/retry} & Xử lý lại tài liệu bị lỗi. \\
\hline
\texttt{DELETE} & \texttt{/v1/rag/documents/\{id\}} & Xóa tài liệu khỏi kho tri thức. \\
\hline
\texttt{GET} & \texttt{/v1/rag/health} & Kiểm tra trạng thái các thành phần RAG. \\
\hline
\end{longtable}

Request truy vấn RAG có thể có dạng:

\begin{verbatim}
POST /api/v1/rag/query

{
  "question": "Điều kiện xét tuyển thẳng của trường là gì?",
  "conversationId": "conv-001",
  "topK": 5,
  "filters": {
    "year": 2025,
    "documentType": "admission_policy"
  }
}
\end{verbatim}

Response RAG cần trả về câu trả lời và nguồn tham chiếu:

\begin{verbatim}
{
  "success": true,
  "data": {
    "answer": "Theo đề án tuyển sinh, điều kiện xét tuyển thẳng gồm ...",
    "sources": [
      {
        "documentName": "De an tuyen sinh 2025.pdf",
        "chunkId": "chunk-001",
        "page": 12,
        "score": 0.87
      }
    ]
  }
}
\end{verbatim}

Thiết kế response có nguồn tham chiếu giúp tăng độ tin cậy của câu trả lời và hỗ trợ người dùng kiểm chứng thông tin. Đồng thời, quản trị viên có thể dựa vào thông tin \texttt{chunkId}, \texttt{documentName} và \texttt{score} để đánh giá chất lượng truy hồi.
